\documentclass[10pt,a4paper]{article}
\usepackage{amsmath,amssymb,amsthm,graphicx}
\usepackage{amsmath,amssymb,amsthm,mathtools}
\usepackage{hyperref}
\usepackage{cleveref}
\newtheorem{theorem}{Theorem}
\newtheorem{fact}[theorem]{Fact}
\newtheorem{definition}[theorem]{Definition}
\newtheorem{lemma}[theorem]{Lemma}
\newtheorem{proposition}[theorem]{Proposition}
\newtheorem{corollary}[theorem]{Corollary}
\newtheorem{claim}[theorem]{Claim}
\newtheorem{remark}[theorem]{Remark}
\newtheorem{observation}[theorem]{Observation}
\DeclareMathOperator*\disc{disc}
\DeclareMathOperator*\err{err}
\DeclareMathOperator*\opt{opt}
\DeclareMathOperator*\Exp{{\bf E}}
\DeclareMathOperator*\Prob{{\bf Pr}}
\newcommand{\polylog}{\operatorname{polylog}}
\let\citep\cite
\newcommand{\bool}{\left\{0,1\right\}}
\newcommand{\boolpm}{\left\{-1,1\right\}}
\newcommand{\sgn}{\mathrm{sgn}}
\newcommand{\ip}{\mathrm{IP}}
\newcommand{\bigo}[1]{O\left(#1\right)}
\newcommand{\numberthis}{\addtocounter{equation}{1}\tag{\theequation}}
\newcommand{\prblm}[1]{{\uppercase{\rm#1}}}
\newcommand{\MCSP}{\prblm{MCSP}}
\newcommand{\MKTP}{\prblm{MKTP}}
\newcommand{\SG}{\prblm{SUM\text{-}GREATER}}
\newcommand{\MKtP}{\mathrm{MKtP}}
\newcommand{\field}[1]{\mathbb{F}_{#1}}
\newcommand{\CNF}{{\sf CNF}}
\newcommand{\DNF}{{\sf DNF}}
\newcommand{\AND}{{\sf AND}}
\newcommand{\OR}{{\sf OR}}
\newcommand{\PARITY}{{\sf PARITY}}
\newcommand{\FORMULA}{{\sf FORMULA}}
\newcommand{\ANY}{{\sf ANY}}
\newcommand{\XOR}{{\sf XOR}}
\newcommand{\LTF}{{\sf LTF}}
\newcommand{\PTF}{{\sf PTF}}
\newcommand{\SYM}{{\sf SYM}}
\newcommand{\GIP}{{\sf GIP}}
\newcommand{\TT}[1]{{\tt tt}\!\left(#1\right)}
\newcommand{\CC}[1]{{\tt CC}\!\left(#1\right)}
\newcommand{\EI}{\aleph}
\newcommand{\given}{\mid}
\newcommand{\uniform}{U}
\newcommand{\LH}[1]{#1^{\mathrm{L}}}
\newcommand{\RH}[1]{#1^{\mathrm{R}}}
\newcommand{\myapprox}[1]{\ensuremath{\stackrel{#1}{\approx}}}
\newcommand{\real}{\mathbb{R}}
\newcommand{\naturals}{\mathbb{N}}
\newcommand{\abs}[1]{\left|#1\right|}
\newcommand{\prn}[1]{\left(#1\right)}
\newcommand{\cprn}[1]{\!\left(#1\right)}
\newcommand{\sqbra}[1]{\left[#1\right]}
\newcommand{\csqbra}[1]{\!\left[#1\right]}
\newcommand{\brkts}[1]{\left\{#1\right\}}
\newcommand{\cbrkts}[1]{\!\left\{#1\right\}}
\newcommand{\boool}{\brkts{-1,1}}
\newcommand{\angles}[1]{\left\langle#1\right\rangle}
\newcommand{\bars}[1]{\left\|#1\right\|}
\newcommand{\cpoly}[1]{\textup{poly}\cprn{#1}}
\newcommand{\quasipoly}[1]{\textup{quasipoly}\prn{#1}}
\newcommand{\cquasipoly}[1]{\textup{quasipoly}\cprn{#1}}
\newcommand{\var}[1]{\textbf{Var}\csqbra{#1}}
\newcommand{\cov}[2]{\textbf{Cov}\csqbra{#1,#2}}
\newcommand{\weight}[2]{{\bf W}^{#2}\csqbra{#1}}
\newcommand{\bin}[1]{{\rm Bin}\prn{#1}}
\newcommand{\cbin}[1]{{\rm Bin}\cprn{#1}}
\newcommand{\pr}[1]{\Prob\sqbra{#1}}
\newcommand{\cpr}[1]{\Prob\csqbra{#1}}
\newcommand{\prt}[2]{\Prob_{#2}\sqbra{#1}}
\newcommand{\cprt}[2]{\Prob_{#2}\csqbra{#1}}
\newcommand{\expect}[1]{\Exp\sqbra{#1}}
\newcommand{\cexpect}[1]{\Exp\csqbra{#1}}
\newcommand{\expectt}[2]{\Exp_{#2}\sqbra{#1}}
\newcommand{\cexpectt}[2]{\Exp_{#2}\csqbra{#1}}
\newcommand{\ACMOD}{{\sf AC}^0\circ{\sf MOD}_2}
\newcommand{\DNFMOD}{{\sf DNF}\circ{\sf MOD}_2}
\newcommand{\ANDMOD}{{\sf AND}\circ{\sf MOD}_2}
\newcommand{\formulaxor}{{\sf FORMULA}\circ{\sf XOR}}
\newcommand{\formulasizexor}[1]{{\sf FORMULA}\csqbra{#1}\circ{\sf XOR}}
\newcommand{\DNFMODAND}{{\sf DNF}\circ{\sf MOD}_2\circ{\sf AND}_t}

% Unused native Set renewal omitted.

\DeclarePairedDelimiter\ceil{\lceil}{\rceil}
\DeclarePairedDelimiter\floor{\lfloor}{\rfloor}
% Duplicate paired-delimiter declaration omitted.
% Duplicate paired-delimiter declaration omitted.
\allowdisplaybreaks
\AtBeginDocument{\author{Valentine Kabanets, Sajin Koroth, Zhenjian Lu\\Dimitrios Myrisiotis, Igor C. Oliveira}}
\newcommand{\AC}{\mathsf{AC}}
\newcommand{\poly}{\operatorname{poly}}
\title{Pseudorandom generators for formulas of low-communication gates via polynomial fooling and NIH recycling}
\author{Valentine Kabanets, Sajin Koroth, Zhenjian Lu, Dimitrios Myrisiotis, Igor C. Oliveira}
\date{doi:10.4230/LIPIcs.CCC.2020.15; CC BY 3.0 Unported}
\begin{document}
\maketitle
\noindent\textit{Editorial note.} Editable transcription of selected formally published CCBY3.0 Unported PRG results; nonexclusive nativev1 private. MainTheorem2 complete construction mechanisms and necessaryAppendixA.2 retained; GIP lowerbound/AppendixA.1 are shared uncounted source-local context, not another task. UnrelatedMCSP sketches removed, other gateclasses/parameters and directMKtP consequence grouped uncounted. Original author parameters/inequalities, including printed extractor bounds in prior-fact statement, preserved without mathematical correction or new assumptions. Wrapper uses correctedSet/ceil/floor dependency handling and independentpreamble copy, metadatafiveauthors unchanged. Noinput/include/assets. All production/validationYupeng; root owns compilation.\par\medskip
\noindent\textit{Allowed standard prerequisites (editorial).}\begin{itemize}
\item Low-degree formula approximation, communication rectangles, explicit extractors and hardness-versus-randomness hybrids. RequiredformalSection2; complete source-local polynomial fooling coefficient argument, NIH recursive construction andAppendixA.2 analysis retained; Section3/GIP proof shared uncounted context
\end{itemize}\medskip
\setcounter{theorem}{1}
	\begin{theorem}[Pseudorandom generators] \label{thm:main_PRG}
    	Let $\mathcal{G}$ be a class of $n$-bits functions. Then,
    	\begin{itemize}
            \item[\emph{1.}] In the context of parity functions, there is a \emph{PRG} that $\varepsilon$-fools $\FORMULA[s] \circ \mathsf{XOR}$ of seed length
            $$
                \ell \; = \; O \! \left (\sqrt{s}\cdot\log(s)\cdot\log(1/\varepsilon) + \log(n)\right).
            $$
            
            \item[\emph{2.}] In the context of two-party randomized communication complexity, there is a \emph{PRG} that $\varepsilon$-fools $\FORMULA[s] \circ \mathcal{G}$ of seed length
            $$
                \ell \; = \;  n/2+O \! \left(\sqrt{s}\cdot\left(R^{(2)}_{\varepsilon/(6s)}(\mathcal{G})+\log(s)\right)\cdot\log(1/\varepsilon)\right).
            $$
            More generally, for every $k(n) \geq 2$, let $\mathcal{G}$ be the class of functions that have $k$-party number-in-hand \emph{(NIH)} $(\varepsilon/6s)$-error randomized communication protocols of cost at most $R^{(k\text{-}\mathsf{NIH})}_{\varepsilon/(6s)}$. There exists a \emph{PRG} that $\varepsilon$-fools $\FORMULA[s]\circ \mathcal{G}$ with seed length
            $$
               \ell \; = \; n/k+O \! \left(\sqrt{s}\cdot \left (R^{(k\text{-}\mathsf{NIH})}_{\varepsilon/(6s)} +\log (s) \right )\cdot \log(1/\varepsilon) +\log(k)\right)\cdot\log(k).
            $$
            
            \item[\emph{3.}] In the setting of $k$-party \emph{NOF} randomized communication complexity, there is a \emph{PRG} that $\varepsilon$-fools $\FORMULA[s] \circ \mathcal{G}$ of seed length
            $$
                \ell \; = \; n-\frac{n}{O \! \left(\sqrt{s}\cdot k \cdot 4^k \cdot \left(  R^{(k)}_{\varepsilon/(2s)}(\mathcal{G})+\log(n)\right)\cdot\log(n/\varepsilon)\right)}.
            $$
    	\end{itemize}
	\end{theorem}

\setcounter{section}{1}
\setcounter{theorem}{7}
	\section{Preliminaries}\label{sec:preliminaries}
	
	\subsection{Notation}
	
	Let $n\in\mathbb{N}$; we denote $\brkts{1,\dots,n}$ by $[n]$, and denote by $\uniform_n$ the uniform distribution over $\bool^n$. We use $\widetilde{O}(\cdot)$ (and $\widetilde{\Omega}(\cdot)$)  to hide polylogarithmic factors. That is, for any $f\colon\mathbb{N}\to\mathbb{N}$, we have that $\widetilde{O}(f(n))=O\cprn{f(n)\cdot\polylog\cprn{f(n)}}$.
	
	In this paper, we will mainly use $\boolpm$ as the Boolean basis. In some parts of this paper, we will use the $\bool$ basis for the simplicity of the presentation. This will be specified in corresponding sections.
	
	\subsection{De Morgan formulas and extensions}
	
	\begin{definition}
	    An $n$-variate \emph{de Morgan formula} is a directed rooted tree; its non-leaf vertices (henceforth, \emph{internal gates}) take labels from $\brkts{{\rm AND},{\rm OR},{\rm NOT}}=\brkts{\land,\lor,\neg}$ and its leaves (henceforth, \emph{variable gates}) take labels from the set of variables $\brkts{x_1,\dots,x_n}$. Each internal gate has bounded in-degree (henceforth, \emph{fan-in}); the {\rm NOT} gate in particular has fan-in $1$ and every variable gate has fan-in $0$. The \emph{size} of a de Morgan formula is the number of its leaf gates.
%
	\end{definition}
	
	In this work, we denote by $\FORMULA[s]$ the class of Boolean functions computable by size-$s$ de Morgan formulas. Let $\mathcal{G}$ denote some class of Boolean functions; then, we denote by $\FORMULA[s]\circ\mathcal{G}$ the class of functions computable by some size-$s$ de Morgan formula where its leaves are labelled by functions in $\mathcal{G}$.
	
%
	
%
%
%
%
%
%
%
%
%
    
%
%
%
%
%
    
    \subsection{Approximating polynomials}
    
    \begin{definition}[Point-wise approximation]
        For a Boolean function $f\colon\boolpm^n\to\boolpm$, we say that the function $\tilde{f}\colon\boolpm^n\to\real$ $\varepsilon$-approximates $f$ if for every $z\in\boolpm^n$,
		\[
			\left| f(z)-\tilde{f}(z)\right| \leq \varepsilon.
		\]
    \end{definition}
    
    We will need the following powerful result for the approximating degree of de Morgan formulas.
    
    \begin{theorem}[\cite{Rei11b}, see also~\cite{ BNRdW07}]\label{thm:approx-polynomial}
		Let $s>0$ be an integer and $0<\varepsilon<1$. Any de Morgan formula $F\colon\boolpm^n\to\boolpm$ of size $s$ has a $\varepsilon$-approximating polynomial of degree $d=O(\sqrt{s}\cdot\log(1/\varepsilon))$. That is, there exists a degree-$d$ polynomial $p\colon\boolpm^n\to\real$ over the reals such that for every $z\in\boolpm^n$,
		\[
			\left| p(z)-F(z)\right| \leq \varepsilon.
		\]
	\end{theorem}
	
	Note that \Cref{thm:approx-polynomial} still holds if we use $\bool$ as the Boolean basis.
    
	\subsection{Communication complexity}
	
	We use standard definitions from communication complexity. In this paper we consider the standard two party model of Yao and its generalizations to multiparty setting. We denote deterministic communication complexity of a Boolean function by $D(f)$ in the two party setting. We refer to \citep{Kushilevitz-Nisan97} for 
	standard definitions from communication complexity.
	
%
%
	    
%
%
	
%
	
	\begin{definition}
	    Let $f\colon\bool^{n}\to\bool$ be a Boolean function. The \emph{communication matrix of $f$}, namely $M_f$, is a $2^{n/2}\times2^{n/2}$ matrix defined by $\prn{M_f}_{x,y}:=f(x,y)$.
	\end{definition}
	
	\begin{definition}
	    A \emph{rectangle} is a set of the form $A\times B$, for $A,B\subseteq\bool^n$. A \emph{monochromatic rectangle} is a rectangle $S$ such that for all pairs $(x,y)\in S$ the value $f(x,y)$ is the same.
    \end{definition}
	
	\begin{lemma}
	    Let $\Pi$ be a protocol that computes $f\colon\bool^{n}\to\bool$ with at most $D$ bits of communication. Then, $\Pi$ induces a partition of $M_f$ into at most $2^D$ monochromatic rectangles.
	\end{lemma}
	
	Given a protocol, its \emph{transcript} is the sequence of bits communicated.
	
	\begin{lemma}
	    For every transcript $z$ of some communication protocol, the set of inputs $(x,y)$ that generate $z$ is a rectangle.
	\end{lemma}
	
%
%
%
	
%
%
%

    Below, we recount the definitions of two multiparty communication models used in this work, namely the number-on-forehead and the number-in-hand models.
	
	\begin{definition}[``Number-on-forehead'' communication model; informal]
	    In the $k$-party ``number-on-forehead'' communication model, 
%
	    there are $k$ players and $k$ strings $x_1,\dots,x_k\in\bool^{n/k}$ and player $i$ gets all the strings except for $x_i$. The players are interested in computing a value $f\cprn{x_1,\dots,x_k}$, where $f\colon\bool^{n}\to\bool$ is some fixed function. 
%
	    We denote by $D^{(k)}(f)$ the number of bits that must be exchanged by the best possible number on forehead protocol solving $f$.
    \end{definition}
        
%
    
    We also use the following weaker communication model.
    
    \begin{definition}[``Number-in-hand'' communication model; informal]
	    In the $k$-party ``number-in-hand'' communication model,
%
	    there are $k$ players and $k$ strings $x_1,\dots,x_k\in\bool^{n/k}$ and player $i$ gets only $x_i$.
%
	    The players are interested in computing a value $f\cprn{x_1,\dots,x_k}$, where $f\colon\bool^{n}\to\bool$ is some fixed function.
%
	    We denote by $D^{(k\text{-}\mathsf{NIH})}(f)$ the number of bits that must be exchanged by the best possible communication protocol.
    \end{definition}
    
    Note that $D^{(k\text{-}\mathsf{NIH})}(f)\leq\prn{1-1/k}\cdot n+1$, for any $n$-variate Boolean function $f$, as if $k-1$ players write on the blackboard their string, then the player that did not reveal her input may compute $f\cprn{x_1,\dots,x_k}$ on her own and then publish it.
    
    For the communication models mentioned above, there are also bounded-error randomized versions, denoted by $R_\delta$, $R_\delta^{(k)}$, and $R^{(k\text{-}\mathsf{NIH})}_\delta$, respectively, where $0<\delta<1$ is an upper bound on the error probability of the protocol. In this setting, the players have access to some shared random string, say $r$, and the aforementioned error probability of the protocol is considered over the possible choices of $r$. Moreover, we require the error to be at most $\delta$ on each fixed choice of inputs.
    
    We can extend the definitions of the communication complexity measures, defined above, to classes of Boolean functions, in a natural way. That is, for any communication complexity measure $M\in\brkts{D,D^{(k)},D^{(k\text{-}\mathsf{NIH})},R_\delta,R_\delta^{(k)},R_\delta^{(k\text{-}\mathsf{NIH})}}$ and for any class of Boolean functions $\mathcal{G}$, we may define
    \[
        M\cprn{\mathcal{G}}:=\max_{g\in\mathcal{G}}M\cprn{g}.
    \]
	We note that throughout this paper, we denote by $n$ the number of input bits for the function regardless the communication models. In the $k$-party communication setting (either NOF or NIH), we assume without loss of generality that $n$ is divisible by $k$.
	

\setcounter{theorem}{16}
	\subsection{Pseudorandomness}
	
	A PRG against a class of functions $\mathfrak{C}$ is a deterministic procedure $G$ mapping short Boolean strings (seeds) to longer Boolean strings, so that $G$'s output ``looks random" to every function in $\mathfrak{C}$.
	
    \begin{definition}[Pseudorandom generators]
    	Let $G\colon\boolpm^{\ell}\to\boolpm^n$ be a function, $\mathfrak{C}$ be a class of Boolean functions, and $0<\varepsilon<1$. We say that $G$ is a \emph{pseudorandom generator of seed length $\ell$ that $\varepsilon$-fools $\mathfrak{C}$} if, for every function $f\in\mathfrak{C}$, it is the case that
    	\[
    	    \abs{\Exp_{z\sim\boolpm^\ell}\csqbra{f\cprn{G\cprn{z}}}-
    	    \Exp_{x\sim\boolpm^n}\csqbra{f\cprn{x}}}\leq\varepsilon.
    	\]
    \end{definition}
	
	A PRG $G$ outputting $n$ bits is called \emph{explicit} if $G$ can be computed in $\poly(n)$ time. All PRGs stated in this paper are explicit.
	

	\paragraph{$\mathsf{GIP}^k_n$.}
	
	The Generalized Inner Product function $\GIP^k_{n}\colon \{0,1\}^{n} \to \{0,1\}$ is defined as
	\[
	    \GIP_{n}^{k}\left(x^{(1)},x^{(2)},\dots,x^{(k)}\right) = \sum_{j=1}^{n/k} \bigwedge_{i=1}^{k} x_{j}^{(i)}~(\mathsf{mod}\;2),
	\]
	where $x^{(i)}\in \{0,1\}^{n/k}$ for each $i\in [k]$.

\setcounter{theorem}{18}
	\section{Lower bounds}\label{sec:lower-bounds}
	
	In this section, we prove an average-case lower bound for the generalized inner product function against $\FORMULA\circ\mathcal{G}$, where $\mathcal{G}$ is the set of functions that have low-cost randomized communication protocols in the number-on-forehead setting. This corresponds to Item 1 of Theorem 1 of the original publication. Items 2 and 3 rely on our PRG constructions, and the proofs are deferred to Section 4 of the original publication.
	
	\begin{theorem}\label{thm:lb-main}
		For any integer $k\geq 2$, $s>0$ and any class of functions $\mathcal{G}$, let $C\colon\boolpm^{n}\to\boolpm$ be a function in $\FORMULA[s]\circ\mathcal{G}$ such that
		\[
		    \Prob_{x\sim\boolpm^{n}}\left[C(x)=\GIP_{n}^{k}(x)\right]\geq 1/2+\varepsilon.
		\]
		Then
		\[
	  		s = \Omega\left( \frac{n^2}{k^2\cdot16^{k}\cdot\left(R^{(k)}_{\varepsilon/(2n^2)}(\mathcal{G})+\log n\right)^2\cdot\log^2(1/\varepsilon)}\right).
		\]
	\end{theorem}
	
	We need a couple useful lemmas from~\cite{Tal16}, whose proofs are presented in \Cref{subsec:formula-lemmas} (\Cref{lem:lb-Tal-1-app} and \Cref{lem:lb-Tal-2-app}) for completeness.
	
	\begin{lemma}[\cite{Tal16}]\label{lem:lb-Tal-1}
	    Let $\mathcal{D}$ be a distribution over $\boolpm^n$, and let $f,C\colon\boolpm^n\to\boolpm$ be such that
	    \[
	        \Prob_{x\sim\mathcal{D}}[C(x)=f(x)]\geq 1/2+\varepsilon.
	    \] 
	    Let $\tilde{C}\colon\boolpm^n\to\real$ be a $\varepsilon$-approximating function of $C$, i.e., for every $x\in\boolpm^n$, $|C(x)-\tilde{C}(x)|\leq \varepsilon$. Then,
	    \[
	        \Exp_{x\sim\mathcal{D}}[\tilde{C}(x)\cdot f(x)]\geq \varepsilon.
	    \]
	\end{lemma}
	
	\begin{lemma}[\cite{Tal16}]\label{lem:lb-Tal-2}
	    Let $\mathcal{D}$ be a distribution over $\boolpm^n$ and let $\mathcal{G}$ be a class of functions. For $f\colon\boolpm^n\to\boolpm$, suppose that $D\colon\boolpm^n\to\boolpm\in \FORMULA[s]\circ\mathcal{G}$ is such that 
	    \[
	        \Prob_{x\sim\mathcal{D}}[D(x)=f(x)]\geq 1/2+\varepsilon_0.
	    \]
	    Then there exists some $h\colon\boolpm^n\to\boolpm\in \XOR_{O\left(\sqrt{s}\cdot\log(1/\varepsilon_0)\right)}\circ \mathcal{G}$ such that
	    \[
	        \Exp_{x\sim\mathcal{D}}[h(x)\cdot f(x)]\geq \frac{1}{s^{O\left(\sqrt{s}\cdot\log(1/\varepsilon_0)\right)}}.
	    \]
	\end{lemma}
    
    We also need the following communication-complexity lower bound for $\GIP$.
    
    \begin{theorem}[{\cite[Theorem 2]{BNS92}}]\label{thm:lb-GIP-bound}
        For any $k\geq 2$, any function that computes $\GIP_{n}^{k}$ on more than $1/2+\delta$ fraction of the inputs (over uniformly random inputs) must have $k$-party deterministic communication complexity at least $\Omega \! \left(n/(k \cdot 4^k) - \log(1/\delta)\right)$.
    \end{theorem}
    
    We first show that if a function correlates with some small formula, whose leaves are functions with low \emph{randomized} communication complexity, then it also correlates non-trivially with some function of relatively low \emph{deterministic} communication complexity.
    
    \begin{lemma}\label{thm:lb-proof-1}
        For any distribution $\mathcal{D}$ over $\boolpm^n$, and any class of functions $\mathcal{G}$, let $f\colon\boolpm^n\to\boolpm$ and $C\colon\boolpm^n\to\boolpm\in \FORMULA[s]\circ\mathcal{G}$ be such that 
	    \[
	        \Prob_{x\sim\mathcal{D}}[C(x)=f(x)]\geq 1/2+\varepsilon.
	    \]
	    Then there exists a function $h$, with $k$-party deterministic communication complexity at most
	    \[
	        O\left(R^{(k)}_{\varepsilon/(2s)}(\mathcal{G})\cdot\sqrt{s}\cdot\log(1/\varepsilon)\right),
	    \]
	    such that
	    \[
	        \Prob_{x\sim\mathcal{D}}[h(x)=f(x)]\geq 1/2+1/s^{O(\sqrt{s}\cdot\log(1/\varepsilon))}.
	    \]
    \end{lemma}
    
    \begin{proof}
        Let $C=F(g_1,g_2\dots,g_s)$ be the function in $\FORMULA[s]\circ\mathcal{G}$, where $F$ is a formula and $g_1,g_2,\dots,g_s$ are leaf functions from the class $\mathcal{G}$. For each  $g_i$, consider a $k$-party randomized protocol $ \Pi_i$ of cost at most $R=R^{(k)}_{\varepsilon/(2s)}(\mathcal{G})$ that has an error $\varepsilon/(2s)$.
        Now consider the following function
        \[
            \tilde{C}(x) \vcentcolon=\Exp_{\Pi_1,\Pi_2,\dots,\Pi_s}\left[D(x)\right],
        \]
        where
        \[
            D(x) \vcentcolon= F(\Pi_1(x),\Pi_2(x),\dots,\Pi_s(x)).
        \]
        Note that for any fixed choice of $(\Pi_1,\Pi_2,\dots,\Pi_s)$, $D$ is a formula whose leaves are functions with \emph{deterministic} communication complexity at most $R$. Next, we show the following.
        
        \begin{claim}\label{claim:lb-approximating}
            The function $\tilde{C}$ $\varepsilon$-approximates $C$.
        \end{claim}
        
        \begin{proof}[Proof of \Cref{claim:lb-approximating}]
            First note that since each $\Pi_i$ is a $(\varepsilon/(2s))$-error randomized protocol, by taking the union bound over the $s$ leaf functions, we have that for every input $x\in\boolpm^n$,
            \[
                \Prob_{\Pi_1,\Pi_2,\dots,\Pi_s}[\Pi_1(x)=g_1(x) \land \Pi_2(x)=g_2(x) \land\dots\land\Pi_s(x)=g_s(x)]\geq 1-\varepsilon/2.
            \]
            Denote by $\mathcal{E}$ the event $\Pi_1(x)=g_1(x) \land \Pi_2(x)=g_2(x)\land\dots\land\Pi_s(x)=g_s(x)$. We have for every $x\in\boolpm^n$,
            \begin{align*}
                \tilde{C}(x) &= \Exp_{\Pi_1,\Pi_2,\dots,\Pi_s}\left[D(x)\right]\\
                &=\Exp\left[D(x)\given \mathcal{E}\right]\cdot\Prob[\mathcal{E}]+\Exp\left[D(x)\given \neg\mathcal{E}\right]\cdot\Prob[\neg\mathcal{E}]\\
                &=C(x)\cdot\Prob[\mathcal{E}]+\Exp\left[D(x)\given \neg\mathcal{E}\right]\cdot\Prob[\neg\mathcal{E}].
            \end{align*}
            On the one hand, we have
            \[
                \tilde{C}(x)=C(x)\cdot\Prob[\mathcal{E}]+\Exp\left[D(x)\given \neg\mathcal{E}\right]\cdot\Prob[\neg\mathcal{E}]\leq  C(x) + \varepsilon/2.
            \]
            On the other hand, we get
            \[
                \tilde{C}(x)=C(x)\cdot\Prob[\mathcal{E}]+\Exp\left[D(x)\given \neg\mathcal{E}\right]\cdot\Prob[\neg\mathcal{E}]\geq C(x)\cdot(1-\varepsilon/2)+(-1)\cdot(\varepsilon/2)\geq C(x)-\varepsilon.
            \]
            This completes the proof of the claim.
        \end{proof}
        
        Now by \Cref{claim:lb-approximating} and \Cref{lem:lb-Tal-1}, we have
        \begin{equation}\label{eq:lb-eq-1}
	        \Exp_{x\sim\mathcal{D}}[\tilde{C}(x)\cdot f(x)]\geq \varepsilon.
        \end{equation}
        By the definition of $\tilde{C}$, \Cref{eq:lb-eq-1} implies that there exists some $D$, which is a formula whose leaves are functions with \emph{deterministic} communication complexity at most $R$, such that
        \[
            \Exp_{x\sim\mathcal{D}}[D(x)\cdot f(x)]\geq \varepsilon,
        \]
        which implies
        \[
            \Prob_{x\sim\mathcal{D}}[D(x)=f(x)]\geq 1/2+ \varepsilon/2.
        \]
        Then by \Cref{lem:lb-Tal-2}, there exists a function $h$, which can be expressed as the $\XOR$ of at most $O(\sqrt{s}\cdot\log(1/\varepsilon))$ leaf functions in $D$, such that
        \[
            \Exp_{x\sim\mathcal{D}}[h(x)\cdot f(x)]\geq \frac{1}{s^{O(\sqrt{s}\cdot\log(1/\varepsilon))}},
        \]
        which again implies
        \[
            \Prob_{x\sim\mathcal{D}}[h(x)=f(x)]\geq \frac{1}{2} + \frac{1}{s^{O\left(\sqrt{s}\cdot\log(1/\varepsilon)\right)}}.
        \]
        Finally, note that the $k$-party deterministic communication complexity of $h$ is at most 
        \[
            O(R\cdot\sqrt{s}\cdot\log(1/\varepsilon)),
        \]
        where $R=R^{(k)}_{\varepsilon/(2s)}(\mathcal{G})$.
    \end{proof}
    
    We are now ready to show \Cref{thm:lb-main}.
    
    \begin{proof}[Proof of \Cref{thm:lb-main}]
        Consider \Cref{thm:lb-proof-1} with $f$ being $\GIP_{n}^{k}$ and $\mathcal{D}$ being the uniform distribution. Consider \Cref{thm:lb-GIP-bound} with $\delta=1/s^{O(\sqrt{s}\cdot\log(1/\varepsilon))}$. We have
        \[
            O\left(R^{(k)}_{\varepsilon/(2s)}(\mathcal{G})\cdot\sqrt{s}\cdot\log(1/\varepsilon)\right)\geq n/(k4^k) - O\left(\sqrt{s}\cdot\log(s)\cdot\log(1/\varepsilon))\right),
        \]
        which implies
        \[
            s\geq \Omega\left( \frac{n^2}{k^2\cdot 16^{k}\cdot\left(R^{(k)}_{\varepsilon/(2n^2)}(\mathcal{G})+\log n\right)^2\cdot\log^2(1/\varepsilon)}\right).\qedhere
        \]
    \end{proof}
	

	\section{Pseudorandom generators}\label{sec:prgs}
	
	Some of our PRGs are obtained from a general framework that allows us to reduce the task of fooling $\FORMULA\circ\mathcal{G}$ to the task of fooling the class of functions which are the parity or conjunction of few functions from $\mathcal{G}$.
	
	\subsection{The general framework}
	
	We show that in order to get a PRG for the class of subquadratic-size formulas with leaf gates in $\mathcal{G}$, it suffices to get a PRG for very simple sublinear-size formulas: either $\XOR\circ\mathcal{G}$ or $\AND\circ\mathcal{G}$.  
	
	\begin{theorem}[PRG for $\FORMULA\circ\mathcal{G}$ from PRG for $\XOR\circ\mathcal{G}$ or $\AND\circ\mathcal{G}$]\label{thm:prg-main}
		Let $\mathcal{G}$ be a class of gates on $n$ bits. For any integer $s>0$ and any $0<\varepsilon<1$, there exists a constant $c>0$ such that the following holds. If a distribution $\mathcal{D}$ over $\boolpm^n$ $\left(2^{-c\cdot\sqrt{s}\cdot\log (s)\cdot \log(1/\varepsilon)}\right)$-fools the $\XOR$ (parity) or the $\AND$ (conjunction) of $c\cdot\sqrt{s}\cdot \log(1/\varepsilon)$ arbitrary functions from $\mathcal{G}$, then $\mathcal{D}$ also $\varepsilon$-fools $\FORMULA[s]\circ\mathcal{G}$.
	\end{theorem}
	
	\begin{proof}\renewcommand{\qedsymbol}{$\square$ (\Cref{thm:prg-main})}
    	We first show the case where $\mathcal{D}$ fools the parity of a few functions from $\mathcal{G}$. The proof can be easily adapted to the case of conjunction.
    	
    	Let $C=F(g_1,g_2\dots,g_s)$ be a function in $\FORMULA[s]\circ\mathcal{G}$, where $F$ is a formula, and $g_1,g_2,\dots,g_s$ are functions from the class $\mathcal{G}$. Let $\uniform$ be the uniform distribution over $\boolpm^n$. 
    	We need to show
    	\begin{equation}\label{eq:prg-1}
    		\Exp[C(\mathcal{D})] \myapprox{\varepsilon} \Exp[C(\uniform)].
    	\end{equation}
    	Let $p$ be a $(\varepsilon/3)$-approximating polynomial for $F$ given by \Cref{thm:approx-polynomial}. Note that the degree of $p$ is
    	\[
    		d=O(\sqrt{s}\cdot\log(1/\varepsilon)).
    	\] 
    	Let us replace $F$, the formula part of $C$, with $p$ and let
    	\[
    		\tilde{C}\vcentcolon=p(g_1,g_2\dots,g_s).
    	\]
    	Since $\tilde{C}$ point-wisely approximates $C$, we have
    	\[
    		\Exp[\tilde{C}(\uniform)] \myapprox{\varepsilon/3} \Exp[C(\uniform)],
    	\]
    	and
    	\[
    		\Exp[\tilde{C}(\mathcal{D})] \myapprox{\varepsilon/3} \Exp[C(\mathcal{D})].
    	\]
    	Then to show \Cref{eq:prg-1}, it suffices to show
    	\[
    		\Exp[\tilde{C}(\mathcal{D})] \myapprox{\varepsilon/3} \Exp[\tilde{C}(\uniform)].
    	\]
    	We have
    	\begin{align*}
        	\label{eq:prg-2}
        	\Exp_{x\sim \mathcal{D}}[\tilde{C}(x)]&=\Exp_{x\sim D}\left[\sum_{\substack{S\subseteq[s]: \\ |S|\leq d}}\hat{p}(S)\cdot\prod_{i\in S}g_i(x)\right]\\
        	&=\sum_{\substack{S\subseteq[s]: \\ |S|\leq d}}\hat{p}(S)\cdot\Exp_{x\sim \mathcal{D}}\left[\prod_{i\in S}g_i(x)\right].
        	\numberthis
    	\end{align*}
    	Now note that for each $S\subseteq[s]$, $\prod_{i\in S}g_i(x)$ computes the $\XOR$ of at most $d$ functions from $\mathcal{G}$.
    	Using the fact the distribution $\mathcal{D}$ $\left(\delta=1/2^{c\cdot\sqrt{s}\cdot\log (s)\cdot \log(1/\varepsilon)}\right)$-fools the $\XOR$ of any $d$ functions from $\mathcal{G}$, we get
    	\begin{align*}
        	\Exp_{x\sim \mathcal{D}}[\tilde{C}(x)]&=\sum_{\substack{S\subseteq[s]: \\ |S|\leq d}}\hat{p}(S)\cdot\Exp_{x\sim D}\left[\prod_{i\in S}g_i(x)\right]\\
        	&=\sum_{\substack{S\subseteq[s]: \\ |S|\leq d}}\hat{p}(S)\cdot\left(\Exp_{x\sim \uniform}\left[\prod_{i\in S}g_i(x)\right]+\delta_S\right) \tag{where $|\delta_S|\leq \delta$} \\
        	&=\sum_{\substack{S\subseteq[s]: \\ |S|\leq d}}\left(\hat{p}(S)\cdot\Exp_{x\sim \uniform}\left[\prod_{i\in S}g_i(x)\right]+\hat{p}(S)\cdot\delta_S\right)\\
        	&=\sum_{\substack{S\subseteq[s]: \\ |S|\leq d}}\hat{p}(S)\cdot\Exp_{x\sim \uniform}\left[\prod_{i\in S}g_i(x)\right]+\sum_{\substack{S\subseteq[s]: \\ |S|\leq d}} \hat{p}(S)\cdot\delta_S\\
        	&=\Exp_{x\sim \uniform}[\tilde{C}(x)]+\sum_{\substack{S\subseteq[s]: \\ |S|\leq d}} \hat{p}(S)\cdot \delta_S.
    	\end{align*}
    	It remains to show
    	\[
    		\left| \sum_{\substack{S\subseteq[s]: \\ |S|\leq d}} \hat{p}(S)\cdot \delta_S \right| \leq \varepsilon/3.
    	\]
    	Note that because $p(z)\in[1-\varepsilon/3,1+\varepsilon/3]$ for every $z\in\boolpm^s$, we have
    	\[
    		|\hat{p}(S)|=\left|\Exp_{z\sim \boolpm^s}\left[p(z)\cdot\prod_{i\in S} z_i\right]\right| \leq 1+\varepsilon/3 < 2.
    	\]
    	Then,
    	\[
    		\left|\sum_{\substack{S\subseteq[s]: \\ |S|\leq d}} \hat{p}(S)\cdot \delta_S \right| \leq  \sum_{\substack{S\subseteq[s]: \\ |S|\leq d}} |\hat{p}(S)|\cdot |\delta_S|\leq \delta \cdot \sum_{\substack{S\subseteq[s]: \\ |S|\leq d}} |\hat{p}(S)| \leq \delta\cdot s^{O(\sqrt{s}\cdot\log(1/\varepsilon))} \leq \varepsilon/3,
    	\]
    	where the last inequality holds for some sufficiently large constant $c$.
    	
    	To show the case of conjunction, we can write the approximating polynomial as the sum of all degree-$d$ monomials, each of which is the $\AND$ of at most $d$ variables. One way to do this is to use the domain $\bool$ instead of $\boolpm$ in the above argument. We need to show that the coefficients in this case still have small magnitude.
    	
    	\begin{claim}\label{claim:prg-magnitude}
    		Let $p\colon\boolpm^{n}\to\real$ be a degree-$d$ polynomial of the form
    		\[
    			p(x)=\sum_{\substack{S\subseteq[n]: \\ |S|\leq d}} \hat{p}(S)\cdot \prod_{i\in S} x_i,
    		\]
    		and let $q\colon\bool^{n}\to\real$ be the corresponding polynomial of $p$ over the domain $\{0,1\}^n$, of the form
    		\[
    			q(y)=\sum_{\substack{T\subseteq[n]: \\ |T|\leq d}} \hat{q}(T)\cdot \prod_{i\in T} y_i.
    		\]
    		Then,
    		\[
    			|q|_1 = \sum_{\substack{T\subseteq[n]: \\ |T|\leq d}} |\hat{q}(T)| \leq n^{O(d)}\cdot \max_{\substack{S\subseteq[n]: \\ |S|\leq d}} |\hat{p}(S)|.
    		\]
    	\end{claim}
    	
    	\begin{proof}\renewcommand{\qedsymbol}{$\square$ (\Cref{claim:prg-magnitude})}
    		We have
    		\begin{align*}
    			q(y_1,y_2,\dots,y_n) &= p(1-2y_1,1-2y_2,\dots,1-2y_n)\\
    			&= \sum_{\substack{S\subseteq[n]: \\ |S|\leq d}} \hat{p}(S)\cdot \prod_{i\in S} (1-2y_i)\\
    			&= \sum_{\substack{S\subseteq[n]: \\ |S|\leq d}} \hat{p}(S)\cdot \left(\sum_{\ell\in \{0,1\}^{|S|}}\prod_{\substack{j\in S:\\ \ell_j=1}}-2y_j \right)\\
    			&= \sum_{\substack{S\subseteq[n]: \\ |S|\leq d}} \sum_{\ell\in \{0,1\}^{|S|}} \hat{p}(S)\cdot(-2)^{|\ell|}\cdot \prod_{\substack{j\in S:\\ \ell_j=1}} y_j. \tag{where $|\ell| = \sum_{i=1}^{|S|}\ell_i$}
    		\end{align*}
    		For a pair $(S,\ell)$ where $S \subseteq[n]$, $|S|\leq d$ and $\ell\in\{0,1\}^{|S|}$, let us define the polynomial $q_{(S,\ell)}$ as
    		\[
    			q_{(S,\ell)}(y)= \hat{p}(S)\cdot(-2)^{|\ell|}\cdot \prod_{\substack{j\in S:\\ \ell_j=1}} y_j.
    		\]
    		Note that there are at most $n^{d}\cdot2^{d}$ many pairs of such $(S,\ell)$'s and for each $(S,\ell)$, we have
    		\begin{align*}
    			|q_{(S,\ell)}|_1 = \left|\hat{p}(S)\cdot(-2)^{|\ell|}\right|\leq 2^{d}\cdot |\hat{p}(S)|.
    		\end{align*}
    	    Finally we have
    		\[
    			|q|_1 = \left|\sum_{(S,\ell)} q_{(S,\ell)}\right|_1\leq \sum_{(S,\ell)} |q_{(S,\ell)}|_1\leq n^{d}\cdot 2^d \cdot 2^d \cdot \max_{\substack{S\subseteq[n]: \\ |S|\leq d}} |\hat{p}(S)|,
    		\]
    		as desired.
    	\end{proof}
    	
    	This completes the proof of \Cref{thm:prg-main}.
	\end{proof}

    \subsection{Formulas of low-communication functions in the number-in-hand setting}
    
    In this subsection, we will use $\bool$ as the Boolean basis.
    
	\begin{theorem}\label{thm:prg-formula-lc}
		For any integers $k\geq 2$, $s>0$ and any $0<\varepsilon<1$, let $\mathcal{G}$ be the class of functions that have $k$-party number-in-hand $(\varepsilon/6s)$-error randomized communication protocols of cost at most $R$. There exists a PRG that $\varepsilon$-fools $\FORMULA[s]\circ \mathcal{G}$ with seed length
		\[
			n/k+O\left(\sqrt{s}\cdot (R +\log (s))\cdot \log(1/\varepsilon) +\log(k)\right)\cdot\log(k).
		\]
	\end{theorem}
    
    We need the following PRG that fools single functions with low communication complexity in the number-in-hand model. The proof is presented in \Cref{appendix:prg_cc} (\Cref{lem:prg-lc-app}) for completeness.

	\begin{theorem}[\cite{ASWZ96,INW94}]\label{thm:prg-lc}
		For any $k\geq 2$, there exists a PRG that $\delta$-fools any $n$-bits functions with $k$-party number-in-hand deterministic communication complexity of at most $D'$, with seed length
		\[
			n/k+O\left(D' + \log(1/\delta)+\log(k)\right)\cdot\log(k).
		\]
	\end{theorem}

    Next, we show a PRG for $\FORMULA\circ\mathcal{G}$, where $\mathcal{G}$ is the class of functions with low-cost communication protocols in the number-in-hand setting. We first show for the case of deterministic protocols.
    
    \begin{theorem}\label{thm:prg-formula-lc-det}
		For any integers $k\geq 2$ and $s>0$, let $\mathcal{G}$ be the class of functions whose $k$-party number-in-hand deterministic communication complexity are at most $D$. There is a PRG that $\varepsilon$-fools $\FORMULA[s]\circ \mathcal{G}$ of size $s$ with seed length
		\[
			n/k+O\left(\sqrt{s}\cdot \log(1/\varepsilon)\cdot (D +\log (s)) +\log(k)\right)\cdot\log(k).
		\]
	\end{theorem}
	
	\begin{proof}
		By \Cref{thm:prg-main}, it suffices to show a PRG that $\left(\delta=1/2^{c\cdot\sqrt{s}\cdot\log (s)\cdot \log(1/\varepsilon)}\right)$-fools every function that is the $\XOR$ of $t=c\cdot\sqrt{s}\cdot \log(1/\varepsilon)$ arbitrary functions from $\mathcal{G}$. Note that such a function has deterministic communication complexity at most $D'=t\cdot D$. Then \Cref{thm:prg-formula-lc-det} follows from \Cref{thm:prg-lc}.
	\end{proof}
	
	We now establish the randomized case.
    
    \begin{proof}[Proof of \Cref{thm:prg-formula-lc}]
    	Let $C$ be a function in $\FORMULA[s]\circ \mathcal{G}$. For each of the leaf functions in $C$, consider a $k$-party number-in-hand randomized protocol of cost at most $R$ that has an error at most $\varepsilon/(6s)$. By taking a union bound over the $s$ leaf functions and by viewing a randomized protocol as a distribution of deterministic protocols (as shown in the proof of \Cref{claim:lb-approximating}), we get the following which is a (point-wisely) $(\varepsilon/3)$-approximating function for $C$:
        \[
            \tilde{C}(x) \vcentcolon=\sum_{i} p_i \cdot D_i(x),
        \]
        where each $p_i\in[0,1]$ is some probability density value (so $\sum_i p_i =1$), and each $D_i$ is a formula whose leaves are functions with \emph{deterministic} communication complexity at most $R$. Then to $\varepsilon$-fool $C$, it suffices to $(\varepsilon/3)$-fool its $(\varepsilon/3)$-approximating function $\tilde{C}$. Also, since $\tilde{C}$ is a convex combination of the $D_i$'s, it suffices to $(\varepsilon/3)$-fools all the $D_i$'s. We will do this using the PRG form \Cref{thm:prg-formula-lc-det}. We get that there exists a PRG that $(\varepsilon/3)$-fools each $D_i$ with seed length
        \[
            n/k+O\left(\sqrt{s}\cdot(R+\log(s))\cdot\log(1/\varepsilon)+\log(k)\right)\cdot\log(k),
        \]
        as desired.
    \end{proof}

%
   
%
   
%
%
%
%
%
%
	
%
	
%
	
%
%
%
%
%
%
	
%
	
%
%
%

%
%
%
%
%
%
%
%
%
%
%
%
%
%
%
%
%
%
%
%
%
%
%
%
%
%
%
%
%
%
%
%
%
%
%
%
%
%
%
%
%
%
%
%
%
%
%
%
%
%
%
%
%
    
%
	
%
    
%
%
%
%
%
%
	
%
%
%
	
%
    
%
%
%
%
%
%
%
%
%
%
%

    \subsection{Applications: Fooling formulas of SYMs, LTFs, XORs, and AC\texorpdfstring{$^0$ circuits}{AC0}}

    \subsubsection{\texorpdfstring{$\FORMULA\circ\SYM$}{FORMULA-SYM} and \texorpdfstring{$\FORMULA\circ\LTF$}{FORMULA-LTF}}
	
	Here, we show how the PRG in \Cref{thm:prg-formula-lc} implies PRGs for $\FORMULA\circ\LTF$ and $\FORMULA\circ\SYM$.
	
	\begin{theorem}\label{thm:prg-formula-sym}
		For any size $s>0$ and $0<\varepsilon<1$, there exists a PRG that $\varepsilon$-fools $\FORMULA[s]\circ \LTF$ with seed length
        \[
			O\left(n^{1/2}\cdot s^{1/4}\cdot \log(n)\cdot \log(n/\varepsilon)\right).
		\]
		For $\FORMULA[s]\circ \SYM$, the seed length is
		\[
		    O\left(n^{1/2}\cdot s^{1/4}\cdot \log(n)\cdot \log(1/\varepsilon) \right).
		\]
	\end{theorem}
	
	We need the fact that the class of $\LTF$ has low communication complexity in the number-in-hand model. Consider the following $k$-party $\SG_m$ problem where the $i$-th party holds a $m$-bit number $z_i$ in hand and they want to determine whether $\sum_{i=1}^k z_i > \theta$, where $\theta$ is a fixed number known to all the parties. Nisan~\cite{Nis94} gave an efficient randomized protocol (with public randomness) for this problem.
    
%
%
%
    
    \begin{theorem}[\cite{Nis94}\footnote{Viola~\cite{Vio15} gave a $\delta$-error randomized protocol for the $k$-party $\SG_m$ problem of cost $O(k\cdot\log(k)\cdot\log(m/\delta))$, which is better than Nisan's protocol when $k=m^{o(1)}$.}]\label{thm:prg-k-party-sum}
	    Let $m>0$ be an integer. For any integer $2\leq k\leq m^{O(1)}$, and any $0<\delta<1$, there exists a $\delta$-error randomized protocol of cost $O(k\cdot\log(m)\cdot\log(m/\delta))$ for the $k$-party $\SG_m$ problem.
	\end{theorem}
	    
    By \Cref{thm:prg-k-party-sum} and the fact that every linear threshold function on $n$ bits has a representation such that the weights are $O(n\log(n))$ integers~\cite{MTT61}, we get the following.
    
    \begin{corollary}\label{cor:prg-k-party-LTF}
         For every $k\geq 2$ and $0<\delta<1$, the $k$-party number-in-hand $\delta$-error randomized communication complexity of $\LTF$ is $O(k\cdot\log(n)\cdot\log(n/\delta))$.
    \end{corollary}
    
	\begin{proof}[Proof of \Cref{thm:prg-formula-sym}]
	    By \Cref{cor:prg-k-party-LTF} and \Cref{thm:prg-formula-lc}, for \emph{every} $k\geq 2$ we get a PRG for $\FORMULA\circ\LTF$ of seed length
        \[
    		n/k+O\left(\sqrt{s}\cdot k\cdot \log(n) \cdot\log(ns/\varepsilon)\cdot\log(1/\varepsilon)+\log(k)\right)\cdot\log(k).
    	\]
	    By choosing 
	    \[
	        k=\frac{n^{1/2}}{s^{1/4}\cdot\log(n)\cdot\log(n/\varepsilon)},
	    \]
	    the claimed seed length follows from a simple calculation.
	
	    For $\FORMULA\circ\SYM$, note that every $n$-bit symmetric function has a deterministic $k$-party number-in-hand communication protocol of cost at most $k\cdot\log(n)$. Then the rest can be shown using a similar argument as above (by choosing $k=n^{1/2}/\left(s^{1/4}\cdot\log(n)\right)$).
	\end{proof}

%
       
%

%
%
%
%
%
%

    \subsubsection{\texorpdfstring{$\FORMULA\circ\XOR$}{FORMULA-XOR}}
    
    For the case of $\FORMULA\circ\XOR$, we get a PRG with better seed length. 
    
	\begin{theorem}\label{thm:prg-formula-xor}
		For any size $s>0$ and $0<\varepsilon<1$, there exists a PRG that $\varepsilon$-fools $\FORMULA[s]\circ \XOR$ with seed length
		\[
			O\left(\sqrt{s}\cdot\log(s)\cdot\log(1/\varepsilon)+\log(n)\right).
		\]
	\end{theorem}
	
	\begin{proof}
	    By \Cref{thm:prg-main}, to fool $\FORMULA[s]\circ\mathcal{G}$, it suffices to $\left(\delta=1/2^{O\left(\sqrt{s}\cdot\log (s)\cdot \log(1/\varepsilon)\right)}\right)$-fool the $\XOR$ of a few functions from $\mathcal{G}$, where $\mathcal{G}$ in this case is the set of all $\XOR$ functions. Note that the $\XOR$ of any set of $\XOR$ functions simply computes some $\XOR$ function. Therefore, we can use small-bias distribution, which fools every $\XOR$ function, to fool $\FORMULA[s]\circ\XOR$. Finally, note that there are known constructions for $\delta$-bias distributions that use $O(\log(n/\delta))$ random bits (see e.g.~\cite{AGHP92}).
	\end{proof}
	
%
%
%
%
%
%
%
%
%
%
    \subsubsection{\texorpdfstring{$\FORMULA\circ\AC^0$}{FORMULA-AC0}}

	Another application of \Cref{thm:prg-main} is to take $\mathcal{G}$ to be the set all functions that can be computed by small constant-depth circuits ($\AC^0$). Note the state-of-the-art PRG against size-$M$ depth-$d$ $\AC^0$ has a seed length of $\log^{d+O(1)}(Mn)\cdot\log(1/\varepsilon)$~\cite{ST19}. Below, let $\AC^0_{d,M}$ denote the class of depth-$d$ circuits of size at most $M$.
	
	\setcounter{theorem}{34}
\begin{theorem}\label{thm:prg-formula-ac0}
		For any size $s,m>0$ and $0<\varepsilon<1$, there exists a PRG that $\varepsilon$-fools $\FORMULA\circ \AC^0_{d,M}$ of size $s$ with seed length
		\[
			\log^{d+O(1)}(Mn)\cdot\sqrt{s}\cdot\log(s)\cdot\log(1/\varepsilon).
		\]
	\end{theorem}
	
%
%
%
%
%
%
%
%
%
	\subsection{Formulas of low number-on-forehead communication leaf gates}
	
	In this section, we show a PRG with mild seed length for formulas of functions with low \emph{multi-party number-on-forehead} communication complexity.
	
	\setcounter{theorem}{36}
\begin{theorem}\label{thm:mild-prg}
	    Let $\mathcal{G}$ be a class of $n$-bits functions. For any size $s>0$, there exists a PRG that $\varepsilon$-fools $\FORMULA[s]\circ \mathcal{G}$, with seed length
		\[
		    n-\frac{n}{O\left(\sqrt{s}\cdot k\cdot4^k\cdot \left( R^{(k)}_{\varepsilon/(2s)}(\mathcal{G})+\log(n)\right)\cdot\log(n/\varepsilon)\right)}.
		\]
	\end{theorem}
	
	The PRG is constructed using the hardness vs. randomness paradigm.
	    
	\subsubsection{Hardness based PRGs}
	
	We show how to construct the PRG using the average-case hardness result for formulas of functions with low multi-party communication complexity (\Cref{thm:lb-main}). We start with some notations. For $x\in\boolpm^m$ and an integer $k$ such that $k$ divides $m$, we consider a partition of $x$ into $k$ equal-sized consecutive blocks and write $x=x^{(1)},x^{(2)},\dots,x^{(k)}$, where $x^{(i)}\in\boolpm^{m/k}$ for each $i\in [k]$.
    	
	\begin{lemma}\label{lem:mild-prg}
	    For any integers $m,t,k>0$ such that $k$ divides $m,t$, let $\mathcal{G}$ be a class of functions on $mt+t$ bits, and let $G\colon\boolpm^{m\times t}\to\boolpm^{mt+t}$ be
		\begin{align*}
			&\quad G(x_1,x_2,\dots,x_t) \\
			&=\left(x_1^{(i)},x_2^{(i)},\dots,x_t^{(i)},\GIP_{m}^{k}\left(x_{(i-1)\cdot(t/k)+1}\right),\GIP_{m}^{k}\left(x_{(i-1)\cdot(t/k)+2}\right),\dots, \GIP_{m}^{k}\left(x_{i\cdot(t/k)+1}\right)\right)_{i\in[k]},
		\end{align*}
		where $x_1,x_2,\dots,x_t\in\boolpm^m$.
		Then $G$ is a PRG that $(t\cdot \varepsilon)$-fools $\FORMULA \circ \mathcal{G}$ of size 
		\[
			s=\Omega\left( \frac{m^2}{k^2\cdot16^{k}\cdot\left(R^{(k)}_{\varepsilon/(2m^2)}(\mathcal{G})+\log m\right)^2\cdot\log^2(1/\varepsilon)}\right).
		\]
	\end{lemma}
	
	\begin{proof}
		The high level idea is as follows. We argue that if there is a $\FORMULA \circ \mathcal{G}$ of the claimed size that breaks the PRG, then there is a $\FORMULA \circ \mathcal{G'}$ of the same size that computes $\GIP$ on $m$ bits, where $\mathcal{G'}$ has a $k$-party communication complexity that is at most that of $\mathcal{G}$ \emph{with respect to the $m$-bit input}, and hence contradicts the $\FORMULA \circ \mathcal{G'}$ complexity of the generalized inner product function. The resulting formula is obtained by fixing some input bits of the original $\FORMULA \circ \mathcal{G}$ which breaks the PRG.

		We use a hybrid argument. First consider the distribution given by $G$, where we replace each $\GIP(x_j)$ ($j\in[t]$) with a uniformly random bit; let us denote those random bits as $\uniform_{j}$ for $j\in[t]$ (note that this is just the uniform distribution). Then for each $j\in[t]$, define $H_j$ to be the distribution that we substitute back $\GIP(x_1),\GIP(x_2),\dots,\GIP(x_{j})$ for the corresponding uniform bits in the previous distribution. 

		For the sake of contradiction, suppose there exists a $\FORMULA \circ \mathcal{G}$ $C$ of size $s$ such that
		\[
			\left| \Prob[C(H_t)=1]-\Prob[F(H_0)=1] \right| > t\cdot\varepsilon.
		\]
		By the triangle inequality, there exists a $1\leq j\leq k$ such that
		\[
			\left| \Prob[C(H_j)=1]-\Prob[C(H_{j-1})=1] \right| >\varepsilon.
		\]
		Then by averaging, there exist some fixings of $x_1,\dots,x_{j-1},x_{j+1},\dots,x_{t}$ and $\uniform_{j+1},\dots,\uniform_{t}$ to $C$ such that the above inequality still holds. Let us denote by $C'$ the circuit obtained by $C$ after such fixings and assume without loss of generality $(k-1)t/k\leq j \leq t$. Then we have
		\begin{equation}\label{eq:distinguish}
		    \left| \Prob\left[C'\left(x_j^{(1)},x_j^{(2)},\dots,x_j^{(k)},\GIP(x_j)\right)=1\right]-\Prob\left[C'\left(x_j^{(1)},x_j^{(2)},\dots,x_j^{(k)},\uniform_j\right)=1\right] \right| > \varepsilon.
		\end{equation}
		By a standard ``unpredictability implies pseudorandomness" argument~\cite{Yao82}, we can show that there is some circuit $C''$, obtained from $C'$ by fixing some value for the last bit, that computes the generalized inner product function on $m$ bits with probability greater than $1/2+\varepsilon$ over uniformly random inputs. Note that the size of $C''$ is the same as $C'$ (hence also $C$) , and also $C''$ can be computed by some $\FORMULA \circ \mathcal{G'}$, where $R^{(k)}_{\delta}(\mathcal{G'})\leq R^{(k)}_{\delta}(\mathcal{G})$ for every $\delta$. This contradicts hardness of $\GIP$ for such circuits (\Cref{thm:lb-main}).
	\end{proof}
	
	We are now ready to prove \Cref{thm:mild-prg}.
	
	\begin{proof}[Proof of \Cref{thm:mild-prg}]
		Consider \Cref{lem:mild-prg}. Let $n=mt+t$, and we have $m=\left(\frac{n}{t}-1\right)$. Then 
		\Cref{lem:mild-prg} gives a PRG that $\varepsilon$-fools $\FORMULA\circ\mathcal{G}$ of size
		\begin{align*}
			s &= \Omega\left( \frac{m^2}{k^2\cdot16^{k}\cdot\left(R^{(k)}_{\varepsilon/(2m^2)}(\mathcal{G})+\log m\right)^2\cdot\log^2(t/\varepsilon)}\right)\\
			&\geq \Omega\left(\left(\frac{n}{t}\right)^2/\left(k^2\cdot16^{k}\cdot\left(R^{(k)}_{\varepsilon/(2n^2)}(\mathcal{G})+\log n\right)^2\cdot\log^2(n/\varepsilon)\right)\right),
		\end{align*}
		which yields
		\[
			t\geq \Omega\left(\frac{n}{\sqrt{s}\cdot k\cdot4^{k}\cdot\left(R^{(k)}_{\varepsilon/(2n^2)}(\mathcal{G})+\log n\right)\cdot\log(n/\varepsilon)}\right).
		\]
		Note that the seed length in this case is $n-t$.
	\end{proof}

	\subsubsection{\texorpdfstring{$\MKtP$}{MKtP} lower bounds}
	
	The PRG in \Cref{thm:mild-prg} is sufficient to give an $\MKtP$ lower bound for formulas of functions with low multi-party communication complexity.
	
	\begin{theorem}\label{thm:prg-mktp}
	    For any integer $s>0$ and any class of $N$-bit function $\mathcal{G}$, if $\MKtP$ on $N$-bit can be computed by some function $\FORMULA[s]\circ \mathcal{G}$, then
	    \[
            s= \frac{N^2}{k^2\cdot16^{k}\cdot R^{(k)}_{1/3}(\mathcal{G})\cdot \polylog(N)}.
	    \]
	\end{theorem}
	
	\begin{proof}
    	Let $C$ be a function in $\FORMULA\circ\mathcal{G}$ of size less than
    	\[
    	    \frac{N^2}{k^2\cdot16^{k}\cdot R^{(k)}_{1/3}(\mathcal{G})\cdot\log^c(N)}
    	\]
    	where $c>0$ is some sufficiently large constant. By \Cref{thm:mild-prg}, we have that there is a PRG that $(1/3)$-fools $C$ and its seed length is
    	\[
             N-\polylog(N).
    	\]
    	Also, since the PRG is polynomial-time computable, we get that for every seed, the output of the PRG has $\mathrm{Kt}$ complexity at most $\theta=N-\polylog(N)$. However, consider the $\MKtP$ function with a threshold parameter $\theta$; this function is not fooled by such a PRG, since it accepts every output of the PRG and rejects a uniformly random string with high probability. 
    \end{proof}
    

\appendix
\setcounter{theorem}{49}
    \section{Proofs of useful lemmas}
    
    \subsection{Useful lemmas for formulas}\label{subsec:formula-lemmas}
    
    The proofs in this section are essentially the same as that of~\cite{Tal16}.
    
    \begin{lemma}[\cite{Tal16}, \Cref{lem:lb-Tal-1} restated]\label{lem:lb-Tal-1-app}
	    Let $\mathcal{D}$ be a distribution over $\boolpm^n$, and let $f,C\colon\boolpm^n\to\boolpm$ be such that
	    \[
	        \Prob_{x\sim\mathcal{D}}[C(x)=f(x)]\geq 1/2+\varepsilon.
	    \] 
	    Let $\tilde{C}\colon\boolpm^n\to\real$ be a $\varepsilon$-approximating function of $C$, i.e., for every $x\in\boolpm^n$, $|C(x)-\tilde{C}(x)|\leq \varepsilon$. Then,
	    \[
	        \Exp_{x\sim\mathcal{D}}[\tilde{C}(x)\cdot f(x)]\geq \varepsilon.
	    \]
	\end{lemma}
	
	\begin{proof}
		Note that since $\tilde{C}$ $\varepsilon$-approximate $C$, we have for every $x\in\boolpm^n$
		\[
			\tilde{C}\cdot C(x)\geq 1-\varepsilon,
		\]
		and 
		\[
			\tilde{C}\cdot (1-C(x))\geq -1-\varepsilon.
		\]
		Then,
		\begin{align*}
			\Exp_{x\sim\mathcal{D}}[\tilde{C}(x)\cdot f(x)]&=\Exp_{x\sim\mathcal{D}}[\tilde{C}(x)\cdot f(x)\given C(x)=f(x)]\cdot\Prob_{x\sim\mathcal{D}}[C(x)=f(x)]\\
			&\qquad+\Exp_{x\sim\mathcal{D}}[\tilde{C}(x)\cdot f(x)\given C(x)\ne f(x)]\cdot\Prob_{x\sim\mathcal{D}}[C(x)\ne f(x)]\\
			&\geq(1-\varepsilon)\cdot \Prob_{x\sim\mathcal{D}}[C(x)=f(x)] + (-1-\varepsilon)\cdot \left(1-\Prob_{x\sim\mathcal{D}}[C(x)=f(x)]\right)\\
			&=2\cdot \Prob_{x\sim\mathcal{D}}[C(x)=f(x)] - 1 -\varepsilon\\
			&\geq 2\cdot (1/2+\varepsilon) - 1 - \varepsilon \geq \varepsilon,
		\end{align*}
		as desired.
	\end{proof}
	
	\begin{lemma}[\cite{Tal16}, \Cref{lem:lb-Tal-2} restated]\label{lem:lb-Tal-2-app}
	    Let $\mathcal{D}$ be a distribution over $\boolpm^n$ and let $\mathcal{G}$ be a class of functions. For $f\colon\boolpm^n\to\boolpm$, suppose that $D\colon\boolpm^n\to\boolpm\in \FORMULA[s]\circ\mathcal{G}$ is such that 
	    \[
	        \Prob_{x\sim\mathcal{D}}[D(x)=f(x)]\geq 1/2+\varepsilon_0.
	    \]
	    Then there exists some $h\colon\boolpm^n\to\boolpm\in \XOR_{O\left(\sqrt{s}\cdot\log(1/\varepsilon_0)\right)}\circ \mathcal{G}$ such that
	    \[
	        \Exp_{x\sim\mathcal{D}}[h(x)\cdot f(x)]\geq \frac{1}{s^{O\left(\sqrt{s}\cdot\log(1/\varepsilon_0)\right)}}.
	    \]
	\end{lemma}
	
    \begin{proof}
		Let 
		\[
			D = F(g_1,g_2\dots,g_s)
		\]
		be a device in $\FORMULA\circ\mathcal{G}$ where $F$ is a formula and $g_1,g_2,\dots,g_s$ are function from $\mathcal{G}$.

		Let $p\colon\boolpm^s\to\real$ be a $\varepsilon_0$-approximating polynomial for $F$ of degree $d=O(\sqrt{s}\cdot\log(1/\varepsilon_0))$. Note that we can write
		\[
			p(z)=\sum_{\substack{S\subseteq[s]: \\ |S|\leq d}}\hat{p}(S)\cdot\prod_{i\in S}z_i.
		\]
		Also, for each $S\subseteq[s]$, we have
		\[
			|\hat{p}(S)|=\left| \Exp_{z\in\boolpm^s}[p(z)\cdot\prod_{i\in S}z_i] \right| \leq  1+\varepsilon_0.
		\]
		Now let
		\[
			\tilde{D} \vcentcolon= p(g_1,g_2\dots,g_s).
		\]
		Note that $\tilde{D}$ is a $\varepsilon_0$-approximating function for $D$. Therefore, by \Cref{lem:lb-Tal-1-app}, we have
		\begin{align*}
			\varepsilon_0 &\leq \Exp_{x\sim\mathcal{D}}[D(x)\cdot f(x)]\\
			&= \Exp_{x\sim\mathcal{D}}\left[\left(\sum_{\substack{S\subseteq[s]: \\ |S|\leq d}}\hat{p}(S)\cdot\prod_{i\in S}g_i \right) \cdot f(x)\right]\\
			&=\sum_{\substack{S\subseteq[s]: \\ |S|\leq d}}\hat{p}(S)\cdot \Exp_{x\sim\mathcal{D}}\left[\prod_{i\in S}g_i\cdot f(x)\right]\\
			&\leq \sum_{\substack{S\subseteq[s]: \\ |S|\leq d}} (1+\varepsilon_0) \cdot \left| \Exp_{x\sim\mathcal{D}}\left[\prod_{i\in S}g_i\cdot f(x)\right] \right|.
		\end{align*}
		The above equation is the sum of at most $s^{O(d)}$ summands. Therefore, there exists some $S\subseteq[s]$ such that
		\[
			\left|\Exp_{x\sim\mathcal{D}}\left[\prod_{i\in S}g_i\cdot f(x)\right]\right| \geq \frac{\varepsilon_0}{(1+\varepsilon_0)\cdot s^{O(d)}}\geq \frac{1}{s^{O\left(\sqrt{s}\cdot\log(1/\varepsilon_0)\right)}},
		\]
		which implies that there exists some $h$, such that either $h=\prod_{i\in S}g_i$ or $h=-\prod_{i\in S}g_i$, and
		\[
			\Exp_{x\sim\mathcal{D}}\left[h(x)\cdot f(x)\right] \geq \frac{1}{s^{O\left(\sqrt{s}\cdot\log(1/\varepsilon_0)\right)}}.
		\]
		Finally, note that such $h$ can be expressed as the $\XOR$ of at most $d$ functions from $\mathcal{G}$.
    \end{proof}
    

    \subsection{PRG for low-communication functions in the number-in-hand setting}\label{appendix:prg_cc}
	
	In this subsection, we show how to fool functions with low communication complexity in the number-in-hand model.
    
	\begin{theorem}[\cite{ASWZ96,INW94}, \Cref{thm:prg-lc} restated]\label{lem:prg-lc-app}
		For any $k\geq 2$, there exists a PRG that $\delta$-fools any $n$-bits functions with $k$-party number-in-hand deterministic communication complexity at most $D'$, with seed length
		\[
			n/k+O\left(D' + \log(1/\delta)+\log(k)\right)\cdot\log(k).
		\]
	\end{theorem}
	
	The PRG in \Cref{thm:prg-lc} is based on the PRG by Impagliazzo, Nisan and Wigderson~\cite{INW94} that is used to derandomize ``network algorithms" and space-bounded computation. We will need to use randomness extractors, which we review below.
	
	\begin{definition}[Min-entropy]
    	Let $X$ be a random variable. The \emph{min-entropy} of $X$, denoted by $H_\infty(X)$, is the largest real number $k$ such that $\Prob\csqbra{X=x}\leq2^{-k}$ for every $x$ in the range of $X$. If $X$ is a distribution over $\boolpm^\EI$ with $H_\infty\cprn{X}\geq k$, then $X$ is called a {\em$\prn{\EI,k}$-source}.
    \end{definition}
    
    \begin{definition}[Extractors]
    	A function $\mathrm{Ext}\colon\boolpm^\EI\times\boolpm^d\to\boolpm^m$ is an \emph{$(k,\varepsilon)$-extractor} if, for any $\prn{\EI,k}$-source $X$, and any test $T\colon\boolpm^m\to\boolpm$, it is the case that
    	\[
    	    \left|\Prob[T(\mathrm{Ext}\cprn{X,\uniform_d}X)=1]-\Prob[T(\uniform_m)=1]\right|\leq \varepsilon.
    	\]
    \end{definition}
	
    \begin{theorem}[{\cite[Theorem 6.22]{Vad12}}]\label{thm:prg-ext}
		For any integer $m,\kappa>0$ and $0<\delta'<0$, there exists an explicit $(\kappa,\delta')$ extractor $\mathrm{Ext}\colon\bool^{m}\times\bool^d\to\bool^{m}$ with $d=O(m-k+\log(1/\delta'))$.
	\end{theorem}

    We are now ready to show \Cref{thm:prg-lc}.
    
	\begin{proof}[Proof of \Cref{thm:prg-lc}]\renewcommand{\qedsymbol}{$\square$ (\Cref{thm:prg-lc})}
	    We first describe the construction of the PRG. In fact, we will construct a sequence of PRGs $G_0,G_1,\dots,G_{\log(k)}$. We begin by specifying the parameters of these PRGs.
	    Let $t=\log(k)$, and let
	    \[
	    	d=O\left(D'+\log(1/\delta)+t\right).
	    \]
	    For $i=0,1,\dots,t$, let
	    \begin{itemize}
	        \item $r_0=n/k$,
	        \item $r_i=r_{i-1}+d$.
	    \end{itemize}
	    Note that we have $r_i=n/k+i\cdot d$. Also, let
	    \[
	    	\mathrm{Ext}_i\colon\bool^{r_i}\times\bool^d\to\bool^{r_i}
	    \]
	    be a $(\kappa_i,\delta')$-extractor from \Cref{thm:prg-ext}, where
	    \[
	    	\kappa_i=r_i-D'-2t - \log(1/\delta)
	    \] 
	    and
	    \[
	        \delta'=\delta/\left(3^t\cdot 2^{D'}\right).
	    \]
	    Note that the seed length of the extractors is $d=O\left(D'+\log(1/\delta)+t\right)$. Finally, define $G_i\colon\bool^{r_i}\to\bool^{n/2^{t-i}}$ recursively as follows
	    \begin{itemize}
	    	\item $G_{0}(a)=a$, where $a\in\bool^{n/k}$.
	    	\item $G_i(a,z)=G_{i-1}(a)\circ G_{i-1}(\mathrm{Ext}_{i-1}(a,z))$, where $a\in\bool^{r_{i-1}}$ and $z\in\bool^d$.
	    \end{itemize}
	    We will show that $G_t\colon\bool^{r_t=n/k+t\cdot d}\to\bool^n$ fools any functions $f$ with $k$-party number-in-hand deterministic communication complexity at most $D'$.
	    First, note that such $f$ can be written as
	    \[
	    	f(x_1,x_2,\dots,x_k)=\sum_{i=1}^{2^{D'}} h^{(i)}_{1}(x_1)\cdot h^{(i)}_{2}(x_2)\cdot\ldots\cdot h^{(i)}_{k}(x_k),
	    \]
	    for some $h^{(i)}_{j}\colon\bool^{n/k}\to \bool$ ($i\in\left[2^{D'}\right], j\in[k]$).
	    Therefore, to show that the PRG $G_t$ $\delta$-fool $f$, it suffices to show that $G_t$ $\left(\delta/2^{D'}\right)$-fools every function $g$ of the form
	    \[
	    	g(x_1,x_2,\dots,x_k)=h_1(x_1)\cdot h_2(x_2)\cdot\ldots\cdot h_k(x_k).
	    \]
	    More specifically we show the following.
	    
	    \begin{claim}\label{lem:prg-lc-claim-1}
	    	For every $k \geq 2$ and $  0\leq i \leq t$, the generator $G_i$ defined above $\left(3^{i}\cdot \delta'\right)$-fools every function $g_i\colon\bool^{n/2^{t-i}}\to\bool$ of the form
	    	\[
	    		g_i(x_1,x_2,\dots,x_{k/2^{t-i}})=h_1(x_1)\cdot h_2(x_2)\cdot\ldots\cdot h_{k/2^{t-i}}(x_{k/2^{t-i}}),
	    	\]
	    	where $x_1,x_2,\dots,x_{k/2^{t-i}}\in\bool^{n/k}$.
	    \end{claim}
	    
	    \begin{proof}\renewcommand{\qedsymbol}{$\square$ (\Cref{lem:prg-lc-claim-1})}
	    	The proof is by induction on $i$. The base case is $i=0$, which is trivial given the definition of $G_{0}$. Now suppose the claim holds for $i-1$, we show the case for $i$. This is done using a hybrid argument. Consider the following four distributions
	    	\begin{itemize}
	    		\item $\mathcal{D}_1=U_{n/2^{t-i}}$,
	    		\item $\mathcal{D}_2=U_{n/2^{t-i+1}}\circ G_{i-1}(U_{r_{i-1}})$,
	    		\item $\mathcal{D}_3=G_{i-1}(U_{r_{i-1}}) \circ G_{i-1}(U'_{r_{i-1}})$ ($U$ and $U'$ are two independent uniform distributions),
	    		\item $\mathcal{D}_4=G_i(U_{r_i})$.
	    	\end{itemize}
	    	We want show show that
	    	\[
	    		\left|\Exp[g_i(\mathcal{D}_1)]-\Exp[g_i(\mathcal{D}_4)]\right|\leq 3^{i}\cdot \delta'.
	    	\]
	    	By the triangle inequality, it suffices to show that
	    	\begin{equation}\label{eq:prg-lc-triangle}
	    		\left|\Exp[g_i(\mathcal{D}_1)]-\Exp[g_i(\mathcal{D}_2)]\right| + \left|\Exp[g_i(\mathcal{D}_2)]-\Exp[g_i(\mathcal{D}_3)]\right| + \left|\Exp[g_i(\mathcal{D}_3)]-\Exp[g_i(\mathcal{D}_4)]\right| \leq  3^{i}\cdot \delta'.
	    	\end{equation}
	    	We show \Cref{eq:prg-lc-triangle} by upper bounding each of the three summands.
	    	
	    	\paragraph{First summand.}
	    	
	    	We show that
            \begin{equation}\label{eq:prg-lc-triangle-1}
				\left|\Exp[g_i(\mathcal{D}_1)]-\Exp[g_i(\mathcal{D}_2)]\right| \leq 3^{i-1}\cdot \delta'.
			\end{equation}
			Let us re-write $g_i$ as
	    	\[
	    		g_i(x_1,x_2,\dots,x_{k/2^{t-i}})=\LH{h}(x_1,x_2,\dots,x_{k/2^{t-i+1}})\cdot \RH{h}(x_{k/2^{t-i+1}+1},x_{k/2^{t-i+1}+2},\dots,x_{k/2^{t-i}}),
	    	\]
	    	where
	    	\[
	    	    \LH{h}(y)\vcentcolon=\prod_{j=1}^{k/2^{t-i+1}}h_i(y) \text{\quad and\quad  } \RH{h}(y)\vcentcolon=\prod_{j=k/2^{t-i+1}}^{k/2^{t-1}}h_i(y).
	    	\]
	    	Then,
	    	\begin{align*}
	    		\Exp[g_i(\mathcal{D}_2)]&=\Exp\left[\LH{h}(U_{n/2^{t-i+1}})\cdot \RH{h}(G_{i-1}(U_{r_{i-1}}))\right]\\
	    		&=\Exp\left[\LH{h}(U_{n/2^{t-i+1}})\right]\cdot\Exp\left[\RH{h}(G_{i-1}(U_{r_{i-1}}))\right]\\
	    		&=\Exp\left[\LH{h}(U_{n/2^{t-i+1}})\right]\cdot\left(\Exp\left[\RH{h}(U_{n/2^{t-i+1}})\right]\pm 3^{i-1}\cdot \delta'\right)
	    		\tag{By the induction hypothesis}\\
	    		&=\Exp\left[\LH{h}(U_{n/2^{t-i+1}})\right]\cdot\Exp\left[\RH{h}(U_{n/2^{t-i+1}})\right] \pm 3^{i-1}\cdot \delta'\\
	    		&=\Exp[g_i(\mathcal{D}_1)]\pm 3^{i-1}\cdot \delta',
	    	\end{align*}
	    	as desired.
	    	
	    	\paragraph{Second summand.}
	    	
	    	By a similar argument, it can be shown that
            \begin{equation}\label{eq:prg-lc-triangle-2}
				\left|\Exp[g_i(\mathcal{D}_2)]-\Exp[g_i(\mathcal{D}_3)]\right| \leq 3^{i-1}\cdot \delta'.
			\end{equation}
			We omit the details here.
	    	
	    	\paragraph{Third summand.}
	    	
	    	We show that
            \begin{equation}\label{eq:prg-lc-triangle-3}
				\left|\Exp[g_i(\mathcal{D}_3)]-\Exp[g_i(\mathcal{D}_4)]\right| \leq \delta'.
			\end{equation}
	    	We have
	    	\begin{align*}
	    		\Exp[g_i(\mathcal{D}_4)]&=\Exp[g_i(G_i(U_{r_i}))]\\
	    		&=\Exp\left[\LH{h}(G_{i-1}(X))\cdot \RH{h}(G_{i-1}(\mathrm{Ext}_{i-1}(X,Z)))\right]
	    		\tag{where $X\sim\bool^{r_{i-1}}$ and $Z\sim\bool^{d}$}\\
	    		&=\Exp[A(X)\cdot B(\mathrm{Ext}_{i-1}(X,Z))]
	    		\tag{where $A(\cdot)= \LH{h}(G_{i-1}(\cdot))$  and $B(\cdot)= \RH{h}(G_{i-1}(\cdot))$}\\
	    		&=\Exp[B(\mathrm{Ext}_{i-1}(X,Z))\given A(X)=1]\cdot\Prob[A(X)=1].
	    	\end{align*}
	    	Similarly, we get
	    	\[
	    		\Exp[g_i(\mathcal{D}_3)]=\Exp[B(U_{r_{i-1}})\given A(X)=1]\cdot\Prob[A(X)=1].
	    	\]
	    	As a result, we have
	    	\begin{align*}\label{eq:prg-lc-triangle-3.1}
	    		&\quad\left|\Exp[g_i(\mathcal{D}_4)]-\Exp[g_i(\mathcal{D}_3)]\right|\\
	    		&=\left|\left(\Exp[B(\mathrm{Ext}_{i-1}(X,Z))\given A(X)=1]-\Exp[B(U_{r_{i-1}})\given A(X)=1]\right)\cdot\Prob[A(X)=1]\right|.
	    		\numberthis
	    	\end{align*}
		    On the one hand, if $\Prob[A(X)=1]\leq \delta'$, then \Cref{eq:prg-lc-triangle-3.1} is at most $\delta'$. On the other hand, if $\Prob[A(X)=1]> \delta'$, then
	        \[
	        	H_{\infty}\cprn{X\given A(X)=1} > r_{i-1}-\log(1/\delta')> r_{i-1}-D'-2t -\log(1/\delta)=\kappa_{i-1}.
	        \]
	        Then by the fact that $\mathrm{Ext}_{i-1}$ is a $(\kappa_{i-1},\delta')$-extractor, we have
	        \[
	            \left|\Exp[B(\mathrm{Ext}_{i-1}(X,Z))\given A(X)=1] - \Exp[B(U_{r_{i-1}})\given A(X)=1]\right| \leq \delta'.
	        \]
	        Therefore, \Cref{eq:prg-lc-triangle-3.1} is at most $\delta'$ and this complete the proof of \Cref{eq:prg-lc-triangle-3}. Finally, note that \Cref{eq:prg-lc-triangle} follows from \Cref{eq:prg-lc-triangle-1}, \Cref{eq:prg-lc-triangle-2} and \Cref{eq:prg-lc-triangle-3}. This completes the proof of \Cref{lem:prg-lc-claim-1}.
	    \end{proof}
	    
	    Given \Cref{lem:prg-lc-claim-1}, \Cref{thm:prg-lc} now follows by letting $i=t$.
	\end{proof}

\begin{thebibliography}{99}
\bibitem{AGHP92}
Noga Alon, Oded Goldreich, Johan H{\aa}stad, and Ren{\'{e}} Peralta.
\newblock Simple construction of almost $k$-wise independent random variables.
\newblock {\em Random Struct. Algorithms}, 3(3):289--304, 1992.
\newblock
  \href{https://doi.org/10.1002/rsa.3240030308}{doi.org/10.1002/rsa.3240030308}.

\bibitem{ASWZ96}
Roy Armoni, Michael~E. Saks, Avi Wigderson, and Shiyu Zhou.
\newblock Discrepancy sets and pseudorandom generators for combinatorial
  rectangles.
\newblock In {\em {FOCS}}, pages 412--421, 1996.
\newblock
  \href{https://doi.org/10.4230/10.1109/SFCS.1996.548500}{doi.org/10.1109/SFCS.1996.548500}.

\bibitem{BNRdW07}
Harry Buhrman, Ilan Newman, Hein R{\"{o}}hrig, and Ronald de~Wolf.
\newblock Robust polynomials and quantum algorithms.
\newblock {\em Theory Comput. Syst.}, 40(4):379--395, 2007.
\newblock
  \href{https://doi.org/10.1007/s00224-006-1313-z}{doi.org/10.1007/s00224-006-1313-z}.

\bibitem{BNS92}
L{\'{a}}szl{\'{o}} Babai, Noam Nisan, and Mario Szegedy.
\newblock Multiparty protocols, pseudorandom generators for logspace, and
  time-space trade-offs.
\newblock {\em J. Comput. Syst. Sci.}, 45(2):204--232, 1992.
\newblock
  \href{https://doi.org/10.1016/0022-0000(92)90047-M}{doi.org/10.1016/0022-0000(92)90047-M}.

\bibitem{INW94}
Russell Impagliazzo, Noam Nisan, and Avi Wigderson.
\newblock Pseudorandomness for network algorithms.
\newblock In {\em {STOC}}, pages 356--364, 1994.
\newblock
  \href{https://doi.org/10.1145/195058.195190}{doi.org/10.1145/195058.195190}.

\bibitem{Kushilevitz-Nisan97}
Eyal Kushilevitz and Noam Nisan.
\newblock {\em Communication Complexity}.
\newblock Cambridge University Press, 1997.

\bibitem{MTT61}
Saburo Muroga, Iwao Toda, and Satoru Takasu.
\newblock Theory of majority decision elements.
\newblock {\em Journal of the Franklin Institute}, 271:376--418, 1961.
\newblock
  \href{https://doi.org/10.1016/0016-0032(61)90702-5}{doi.org/10.1016/0016-0032(61)90702-5}.

\bibitem{Nis94}
Noam Nisan.
\newblock The communication complexity of threshold gates.
\newblock In {\em Proceedings of ``Combinatorics, Paul Erdos is Eighty"}, pages
  301--315, 1994.

\bibitem{Rei11b}
Ben Reichardt.
\newblock Reflections for quantum query algorithms.
\newblock In {\em {SODA}}, pages 560--569, 2011.
\newblock \href{https://arxiv.org/abs/1005.1601}{arXiv:1005.1601}.

\bibitem{ST19}
Rocco~A. Servedio and Li-Yang Tan.
\newblock Improved pseudorandom generators from pseudorandom multi-switching
  lemmas.
\newblock In {\em {APPROX/RANDOM}}, pages 45:1--45:23, 2019.
\newblock
  \href{https://doi.org/10.4230/LIPIcs.APPROX-RANDOM.2019.45}{doi.org/10.4230/LIPIcs.APPROX-RANDOM.2019.45}.

\bibitem{Tal16}
Avishay Tal.
\newblock The bipartite formula complexity of inner-product is quadratic.
\newblock {\em Electronic Colloquium on Computational Complexity {(ECCC)}},
  23:181, 2016.
\newblock \href{http://eccc.hpi-web.de/report/2016/181}{ECCC:TR16-181}.

\bibitem{Vad12}
Salil~P. Vadhan.
\newblock Pseudorandomness.
\newblock {\em Foundations and Trends in Theoretical Computer Science},
  7(1-3):1--336, 2012.
\newblock
  \href{https://doi.org/10.1561/0400000010}{doi.org/10.1561/0400000010}.

\bibitem{Vio15}
Emanuele Viola.
\newblock The communication complexity of addition.
\newblock {\em Combinatorica}, 35(6):703--747, 2015.
\newblock
  \href{https://doi.org/10.1145/10.1007/s00493-014-3078-3}{doi.org/10.1007/s00493-014-3078-3}.

\bibitem{Yao82}
Andrew~Chi{-}Chih Yao.
\newblock Theory and applications of trapdoor functions.
\newblock In {\em {FOCS}}, pages 80--91, 1982.
\newblock
  \href{https://doi.org/10.1109/SFCS.1982.45}{doi.org/10.1109/SFCS.1982.45}.

\end{thebibliography}
\end{document}
